\documentclass[10pt]{article}

\usepackage[margin=0.35in]{geometry}
\usepackage{multicol}
\usepackage{amsmath,amssymb,amsfonts}

\makeatletter
\renewcommand\section{\@startsection{section}{1}{0pt}{0.8ex}{0.4ex}{\normalfont\normalsize\bfseries}}
\renewcommand\subsection{\@startsection{subsection}{2}{0pt}{0.5ex}{0.2ex}{\normalfont\bfseries}}
\def\@listi{\leftmargin\leftmargini \topsep 0pt \parsep 0pt \itemsep 0pt}
\makeatother

\setlength{\parindent}{0pt}
\setlength{\parskip}{0.1em}
\setlength{\columnsep}{0.15in}
\setlength{\abovedisplayskip}{3pt}
\setlength{\belowdisplayskip}{3pt}
\setlength{\abovedisplayshortskip}{2pt}
\setlength{\belowdisplayshortskip}{2pt}
\pagestyle{empty}

\begin{document}

{\Large \textbf{SEG800 DIP Cheatsheet}}\\
{\small Test 4: Spatial Filtering, Frequency-Domain Filtering, Color Fundamentals}

\vspace{0.25em}
\hrule
\vspace{0.5em}

\scriptsize
\begin{multicols}{2}

\section{Part I: Spatial Filtering}

\textbf{Noise model:}
\[
\hat I(i,j)=I(i,j)+n(i,j)
\]
Gaussian noise: small random variations. Salt-and-pepper: extreme dark/bright spikes.

\subsection{Linear Filtering}
\[
g(x,y)=\sum_{s=-a}^{a}\sum_{t=-b}^{b}w(s,t)f(x+s,y+t)
\]
Kernel normalization for smoothing:
\[
\sum_s\sum_t w(s,t)=1
\]

\subsection{Correlation vs Convolution}
\[
\text{corr: } (w\circledast f)(x,y)=\sum_s\sum_t w(s,t)f(x+s,y+t)
\]
\[
\text{conv: } (w\star f)(x,y)=\sum_s\sum_t w(s,t)f(x-s,y-t)
\]
Convolution flips the kernel; correlation does not. Symmetric kernels give the same output. Same-size odd $m\times n$ filtering: zero-pad $\lfloor m/2\rfloor,\lfloor n/2\rfloor$ each side; for $3\times 3$, pad 1.

\subsection{Core Kernels}
Box:
\[
\frac{1}{9}
\begin{bmatrix}
1&1&1\\
1&1&1\\
1&1&1
\end{bmatrix}
\qquad
\frac{1}{25}
\begin{bmatrix}
1&1&1&1&1\\
1&1&1&1&1\\
1&1&1&1&1\\
1&1&1&1&1\\
1&1&1&1&1
\end{bmatrix}
\]
Bilinear:
\[
\frac{1}{16}
\begin{bmatrix}
1&2&1\\
2&4&2\\
1&2&1
\end{bmatrix}
\]
Gaussian:
\[
G(x,y)=\exp\!\left(-\frac{x^2+y^2}{2\sigma^2}\right)
\]
Worksheet Q4 template with $\sigma=2$ and $3\times 3$ ($x,y\in\{-1,0,1\}$):
First compute the \textbf{unnormalized} entries:
\[
G(-1,-1)=e^{-2/8}=0.7788,\quad G(-1,0)=e^{-1/8}=0.8825,\quad G(0,0)=1
\]
Use symmetry for the rest, so
\[
G\approx
\begin{bmatrix}
0.7788&0.8825&0.7788\\
0.8825&1.0000&0.8825\\
0.7788&0.8825&0.7788
\end{bmatrix},
\qquad
\sum G\approx 7.6452
\]
Then \textbf{normalize} by dividing every entry by the sum:
\[
h(x,y)=\frac{G(x,y)}{\sum_{x,y}G(x,y)}
\]
\[
h=\frac{1}{7.6452}
\begin{bmatrix}
0.7788&0.8825&0.7788\\
0.8825&1.0000&0.8825\\
0.7788&0.8825&0.7788
\end{bmatrix}
\approx
\begin{bmatrix}
0.1019&0.1154&0.1019\\
0.1154&0.1308&0.1154\\
0.1019&0.1154&0.1019
\end{bmatrix}
\]

\subsection{Median Filter}
For a $3\times 3$ median filter: collect 9 values, sort, choose the 5th. Best for salt-and-pepper noise; usually preserves edges better than averaging.

\subsection{Separable Kernels}
If
\[
w(x,y)=w_1(x)w_2(y)
\]
then filter in two 1-D passes. Proof idea: $(w_1w_2)*f=w_1*(w_2*f)=w_2*(w_1*f)$.
Cost comparison for image $M\times N$, kernel $m\times n$:
\[
\text{direct}=MNmn, \qquad \text{separable}=MN(m+n)
\]
Speed-up:
\[
\frac{MNmn}{MN(m+n)}=\frac{mn}{m+n}
\]
Square $m\times m$ case:
\[
O(MNm^2) \to O(MNm)
\]
For $5\times 5$: $25\to 10$ ops/pixel, speed-up $=2.5\times$.
Worksheet factorization example:
\[
w=
\begin{bmatrix}
1&3&1\\
2&6&2
\end{bmatrix}
\]
Pick $(2,1)$, so $E=2$,
\[
\mathbf c=
\begin{bmatrix}
1\\2
\end{bmatrix}
\!,\ \mathbf r=
\begin{bmatrix}
2&6&2
\end{bmatrix}
\!,\ w_1=\mathbf c,
\ w_2=\frac{\mathbf r}{E}=
\begin{bmatrix}
1&3&1
\end{bmatrix}
\]
Check:
\[
w_1w_2=
\begin{bmatrix}
1\\2
\end{bmatrix}
\begin{bmatrix}
1&3&1
\end{bmatrix}=
\begin{bmatrix}
1&3&1\\2&6&2
\end{bmatrix}
\]

\subsection{Gaussian Pixel Template}
For a target pixel, use its local neighborhood and compute
\[
g(i,j)=\sum_x\sum_y h(x,y)f(i+x,j+y)
\]
Worksheet image at $(2,3)$ uses neighborhood
\[
\begin{bmatrix}
100&102&98\\
101&99&103\\
97&100&102
\end{bmatrix}
\]
so
\[
\begin{aligned}
g(2,3)\approx{}&100(0.1019)+102(0.1154)+98(0.1019)\\
&+101(0.1154)+99(0.1308)+103(0.1154)\\
&+97(0.1019)+100(0.1154)+102(0.1019)
\end{aligned}
\]
\[
g(2,3)\approx 100.26\approx 100
\]
For high-boost, first compute $\bar f(i,j)$ this way, then use $g_{\text{mask}}=f-\bar f$.

\subsection{Sharpening}
Laplacian:
\[
\nabla^2 f=f(x+1,y)+f(x-1,y)+f(x,y+1)+f(x,y-1)-4f(x,y)
\]
Common kernel:
\[
\begin{bmatrix}
0&-1&0\\
-1&4&-1\\
0&-1&0
\end{bmatrix}
\]
Sharpening:
\[
g=f+c\nabla^2 f
\]
Choose sign of $c$ to match the kernel convention.

Unsharp mask and high-boost:
\[
g_{\text{mask}}=f-\bar f
\]
\[
g=f+k g_{\text{mask}}=(1+k)f-k\bar f
\]
$k=1$: unsharp masking. $k>1$: high-boost.

\section{Part II: Frequency-Domain Filtering}

\subsection{2-D DFT / IDFT}
\[
F(u,v)=\sum_{x=0}^{M-1}\sum_{y=0}^{N-1}f(x,y)e^{-j2\pi(ux/M+vy/N)}
\]
\[
f(x,y)=\frac{1}{MN}\sum_{u=0}^{M-1}\sum_{v=0}^{N-1}F(u,v)e^{j2\pi(ux/M+vy/N)}
\]
Magnitude, phase, power:
\[
F=R+jI=|F|e^{j\phi}
\]
\[
|F|=\sqrt{R^2+I^2}, \qquad \phi=\operatorname{atan2}(I,R), \qquad P=|F|^2
\]
DC component:
\[
F(0,0)=MN\bar f
\]

\subsection{Sampling and Aliasing}
1-D Nyquist:
\[
f_s>2f_{\max}
\]
2-D Nyquist:
\[
\frac{1}{\Delta T}>2\mu_{\max}, \qquad \frac{1}{\Delta Z}>2\nu_{\max}
\]
Undersampling causes aliasing and moire patterns.

\subsection{Translation and Centering}
Spatial shift changes phase, not magnitude:
\[
f(x-x_0,y-y_0) \Longleftrightarrow F(u,v)e^{-j2\pi(ux_0/M+vy_0/N)}
\]
Center spectrum before filtering:
\[
f(x,y)(-1)^{x+y} \Longleftrightarrow F\left(u-\frac{M}{2},v-\frac{N}{2}\right)
\]

\subsection{Convolution Theorem and Padding}
\[
f*h \Longleftrightarrow FH
\]
DFT multiplication gives circular convolution unless padded enough.
For $A\times B$ convolved with $C\times D$:
\[
P\ge A+C-1, \qquad Q\ge B+D-1
\]
Equal-size shortcut:
\[
P\ge 2M-1, \qquad Q\ge 2N-1
\]
Common convenient choice:
\[
P=2M, \qquad Q=2N
\]
Butterworth/other radial filters are built on the padded grid $P\times Q$, centered at $(P/2,Q/2)$. In this course, if not stated otherwise, usually use $P=2M$, $Q=2N$ before computing $D(u,v)$.

\subsection{Frequency-Domain Filtering Workflow}
\begin{enumerate}
\item Pad image to $P\times Q$.
\item Multiply by $(-1)^{x+y}$.
\item Compute DFT $F(u,v)$.
\item Build centered filter $H(u,v)$.
\item Multiply: $G(u,v)=H(u,v)F(u,v)$.
\item IDFT, multiply by $(-1)^{x+y}$, take real part, crop.
\end{enumerate}

\subsection{Radial Distance and Filter Families}
\[
D(u,v)=\sqrt{\left(u-\frac{P}{2}\right)^2+\left(v-\frac{Q}{2}\right)^2}
\]
Ideal LPF:
\[
H_{ILPF}=
\begin{cases}
1,&D\le D_0\\
0,&D>D_0
\end{cases}
\]
Gaussian LPF:
\[
H_{GLPF}=e^{-D^2/(2D_0^2)}
\]
Butterworth LPF:
\[
H_{BLPF}=\frac{1}{1+(D/D_0)^{2n}}
\]
Worksheet template: find padded center $\to$ compute $D(u,v)$ $\to$ substitute into $H_{BLPF}$.
High-pass from low-pass:
\[
H_{HP}=1-H_{LP}
\]
Gaussian HPF:
\[
H_{GHPF}=1-e^{-D^2/(2D_0^2)}
\]
Butterworth HPF:
\[
H_{BHPF}=\frac{1}{1+(D_0/D)^{2n}}, \qquad H_{BHPF}(0,0)=0
\]

\subsection{Behavior and Computation}
Small $D_0$ in LPF $\Rightarrow$ stronger smoothing. Larger $D_0$ in HPF $\Rightarrow$ stronger removal of low frequencies.

Ideal filters have strongest ringing. Gaussian has least ringing. Butterworth is a tunable compromise via order $n$.

Ops estimate:
\[
\text{spatial}=MNmn, \qquad \text{FFT}\approx 2MN\log_2(MN)
\]
Use FFT when kernels are large and nonseparable.

\section{Part III: Color Fundamentals}

\subsection{Terms and Primaries}
Radiance = physical energy, luminance = perceived light, brightness = subjective, intensity = numeric proxy.

Additive light primaries: RGB.
\[
R+G=Y, \qquad G+B=C, \qquad R+B=M, \qquad R+G+B=W
\]
Subtractive primaries: CMY. Printing uses CMYK.

\subsection{Chromaticity and Gamut}
\[
x=\frac{X}{X+Y+Z}, \qquad y=\frac{Y}{X+Y+Z}, \qquad z=\frac{Z}{X+Y+Z}, \qquad x+y+z=1
\]
Gamut = set of colors a device can reproduce.

\subsection{RGB, CMY, CMYK}
Normalized RGB cube:
\[
0\le R,G,B\le 1, \qquad R=G=B \Rightarrow \text{gray}
\]
\[
(0,0,0)=\text{black}, \qquad (1,1,1)=\text{white}
\]
24-bit RGB gives $256^3=2^{24}=16{,}777{,}216$ colors.

RGB $\leftrightarrow$ CMY:
\[
C=1-R, \qquad M=1-G, \qquad Y=1-B
\]
\[
R=1-C, \qquad G=1-M, \qquad B=1-Y
\]

CMYK from CMY:
\[
K=\min(C,M,Y)
\]
If $K\ne 1$:
\[
C'=\frac{C-K}{1-K}, \qquad M'=\frac{M-K}{1-K}, \qquad Y'=\frac{Y-K}{1-K}
\]

\subsection{HSI}
Hue = basic color, saturation = purity, intensity = brightness.
\[
I=\frac{R+G+B}{3}
\]
\[
S=1-\frac{3}{R+G+B}\min(R,G,B)
\]
\[
\theta=\cos^{-1}\!\left(\frac{\frac12[(R-G)+(R-B)]}{\sqrt{(R-G)^2+(R-B)(G-B)}}\right)
\]
\[
H=
\begin{cases}
\theta,& B\le G\\
360^\circ-\theta,& B>G
\end{cases}
\]
Hue sectors:
\[
0^\circ\le H<120^\circ: RG, \quad 120^\circ\le H<240^\circ: GB, \quad 240^\circ\le H<360^\circ: BR
\]

\subsection{Color Processing Rules}
Two main approaches:
\begin{itemize}
\item Process each RGB channel and recombine.
\item Convert to HSI and process only $I$ when the task is brightness, smoothing, sharpening, or histogram equalization.
\end{itemize}
Do not histogram-equalize RGB channels independently: different nonlinear mappings can distort hue.

Brightness scaling:
\[
R'=kR, \qquad G'=kG, \qquad B'=kB
\]
or in HSI:
\[
I'=kI
\]

RGB complement:
\[
R'=1-R, \qquad G'=1-G, \qquad B'=1-B
\]

\subsection{Pseudocolor and Slicing}
Intensity slicing:
\[
\mathbf g(x,y)=\mathbf c_k \quad \text{if } l_{k-1}\le f(x,y)<l_k
\]
Intensity-to-color mapping:
\[
R=T_R(f), \qquad G=T_G(f), \qquad B=T_B(f)
\]

Color distance in RGB:
\[
D=\sqrt{(R-a_R)^2+(G-a_G)^2+(B-a_B)^2}
\]
for color slicing around a target color.

\end{multicols}
\end{document}
